\documentclass[11pt]{article}
\usepackage[margin=1in]{geometry}
\usepackage{times,amsmath,amssymb,xcolor,enumitem,microtype,parskip}
\usepackage[hidelinks]{hyperref}
\setlength{\parindent}{0pt}
\definecolor{reviewblue}{RGB}{0,74,135}
\newenvironment{reviewcomment}{\begin{quote}\color{reviewblue}\itshape}{\end{quote}}
\newcommand{\Response}{\par\noindent\textbf{Response.}\ }
\newcommand{\Changes}{\par\noindent\textbf{Changes in the revised manuscript.}\ }
\newcommand{\Location}[1]{\par\noindent\textit{Location in the revised manuscript: #1}}
\begin{document}
\begin{center}
{\LARGE\bfseries Response Letter}\\[0.8em]
{\large Manuscript ID: 26-3482}\\[0.35em]
{\large ADAPT: Adaptive Dual-projection Architecture for Perceptive Traversal}
\end{center}
\vspace{0.8em}
Dear Associate Editor and Reviewers,

Thank you for your careful assessment and constructive suggestions. We have revised the manuscript to clarify the efficiency comparisons and adaptive sensing mechanism, explain the adaptive reward, and improve the presentation of the simulation and hardware experiments.

The principal revisions are summarized below:
\begin{itemize}
    \item We have added deployment latency and resource-usage benchmarks against Gallant on the same local computer (Table~IV(b), Sec.~IV-C-3), showing $29.6\%$ lower mean perception-to-command-publication latency and $58.4\%$ lower sampled peak GPU memory usage, together with lower CPU utilization and RAM usage.
    \item We have specified the Gallant comparison: $98.6\%$ fewer per-frame pre-encoder exteroceptive values, approximately $56\%$ lower peak GPU memory, and $32.7\%$ of Gallant's training-iteration time under the stated training settings.
    \item We have clarified that adaptive perception uses radius smoothness and regularization rather than a speed- or obstacle-dependent target radius, and specified the reward formulas and weights in the detailed responses.
    \item We have clarified that the adaptive radius changes only the vertical distance map; the horizontal elevation map retains fixed coverage and sampling resolution.
    \item We have specified that the adaptive-radius ablation in Table~V is a simulation experiment, identified which two reward terms are removed, and removed the Table~V comparison from Sec.~IV-C-1.
    \item We have revised Fig.~1 and Sec.~III-C to distinguish the actor from the privileged critic and show separate recurrent LSTM states. We have added error bars showing the sample standard deviation across ten independent training seeds in Fig.~5 (Sec.~IV-C-2), and corrected the typo in Sec.~I.
\end{itemize}

In the detailed response below, the Associate Editor's points are summarized and the reviewers' comments are reproduced with minor typographic normalization in \textcolor{reviewblue}{blue italics}; our responses appear in black. Each response identifies the relevant manuscript sections and figures or tables; clarifications provided only in this letter are explicitly distinguished from manuscript changes.

\clearpage
\section*{Response to the Associate Editor's Report (ReportOn26-3482v1)}
\subsection*{AE Overall assessment}
\begin{reviewcomment}
Both reviewers find the problem relevant and well motivated, and recognize the interest of adjusting how far ahead the robot perceives obstacles for different 3D locomotion tasks. The report summarizes ADAPT's horizontal and vertical projections, policy-controlled sensing radius, simulation comparisons with PIM, Humanoid Parkour, TSDF, and Gallant, and zero-shot deployment on a Unitree G1.
\end{reviewcomment}
\Response We sincerely thank the Associate Editor for the thoughtful assessment and for recognizing the motivation and practical relevance of the work. We also appreciate the careful summary of the dual-projection representation, adaptive sensing radius, baseline evaluations, and real-world deployment. The specific concerns raised in the report have helped us clarify the scope and evidence of our efficiency claims and improve the presentation of the method and experiments.

\subsection*{AE Comment 1: Evidence for compactness and computational efficiency}
\begin{reviewcomment}
Authors claim that ADAPT provides a more compact representation and potentially makes reinforcement learning more computationally efficient, compared with voxel-based alternatives. The authors should specify how the comparisons were conducted. They have to provide enough detail to verify both the dimensionality reduction and the training-speed comparison. In particular, R8 asks where the 56\% computational overhead reduction and the 98.6\% observation-dimensionality reduction are supported. R3 notes that Table IV reports only training efficiency, whereas inference efficiency matters more for real-time perceptive control. R3 also asks for latency and onboard resource measurements against the main baselines.
\end{reviewcomment}
\Response Thank you for requesting this clarification. We agree that the efficiency metrics should be stated explicitly. The 56\% figure refers to \emph{peak GPU memory during training} relative to the Gallant voxel baseline. Using the same modified NVIDIA RTX 4090 GPU with 48 GB of VRAM and identical PPO settings (4,096 parallel environments, a rollout of 24 steps per environment per policy update, 5 PPO epochs per update, and 4 mini-batches per epoch; each rollout batch contains $4{,}096\times24=98{,}304$ transitions), peak memory is 36,054\,MiB for Gallant and 15,870\,MiB for ADAPT: $(36{,}054-15{,}870)/36{,}054=55.98\%$. The respective iteration times are 34.6\,s and 11.3\,s, so ADAPT takes 32.7\% of Gallant's time under this setup.

Sec.~IV-B-3 states that the methods share the same training settings; the specific numerical PPO settings listed above are provided only in this response letter.

The 98.6\% figure compares \emph{single-frame external perception inputs before encoding}. Gallant's $32\times32\times40$ voxel grid has 40,960 values; ADAPT's $21\times17$ horizontal elevation and $13\times17$ vertical distance maps have 578 values in total. Thus $(40{,}960-578)/40{,}960=98.59\%$.

To evaluate deployment efficiency, we added measurements of deployment latency and computational resource usage. Table~IV(b) compares ADAPT with the main voxel baseline, Gallant, on the same local computer (Intel i7-12700H CPU and NVIDIA RTX 3060 Laptop GPU). This local computer was also used for perception and policy inference in the real-world robot experiments; we have corrected the deployment description in Sec.~IV-A-2. Mean map-construction latency decreases from 64.06 to 41.64\,ms, policy inference latency on GPU decreases from 1.91 to 1.40\,ms, and perception-to-command-publication latency from 79.45 to 55.94\,ms (29.6\% lower). Sampled peak GPU memory decreases from 5,711 to 2,373\,MiB (58.4\% lower); mean whole-host CPU utilization decreases from 38.19\% to 17.76\%, and RAM usage from 10.66 to 8.04\,GiB.
\Changes
\begin{itemize}
    \item Sec.~I restates the two efficiency quantities.
    \item Table~I in Sec.~II-A defines the pre-encoder sensory dimension; Sec.~IV-B-3 derives the comparison from these dimensions and the training measurements in Table~IV(a).
    \item Table~IV(b) and Sec.~IV-C-3 report the local-computer runtime comparison, measurement boundaries, hardware, and trial setup.
\end{itemize}
\Location{Abstract and Sec.~I, p.~1; Sec.~II-A and Table~I, p.~2; Sec.~IV-B-3, p.~5; Table~IV(a)--(b), p.~6; Sec.~IV-C-3, p.~7.}



\subsection*{AE Comment 2: Voxel baselines and adaptive reward}
\begin{reviewcomment}
Voxel-based representations are not inherently inferior and can have different computational trade-offs. Analyze whether the voxel-based baselines' performance gap arises from optimization difficulty, network capacity, training budget, representation redundancy, or implementation factors. Explain the role of the adaptive reward in learning the radius and clarify whether it encodes a hand-designed heuristic.
\end{reviewcomment}
\Response Thank you for this helpful suggestion. We agree that the baseline results warrant further analysis. We have added an analysis in Sec.~IV-B-3 (Performance Comparison). Voxel representations retain richer 3D geometry and are not inherently inferior. For the Gallant voxel baseline implemented in this study, however, fine spatial resolution limits sensing coverage under a fixed voxel count and computational budget. The success-rate gap relative to ADAPT is smaller at lower speeds and widens at high speed. In high-speed tests, we observed abrupt deceleration after obstacles entered the sensing range, sometimes followed by loss of balance. We attribute this behavior to limited look-ahead, which may leave insufficient time for smooth adjustments of speed and posture and thus increase the difficulty of high-speed traversal. In contrast, ADAPT allows the policy to adjust its sensing range, and hence its look-ahead distance, according to the robot's motion state.

The adaptive reward contains radius smoothness and regularization terms. With current radius $r_t$, previous radius $r_{t-1}$, and initial radius $r_{\mathrm{init}}$, their contributions are $R_{\mathrm{smooth}}(t)=-\lambda_s|r_t-r_{t-1}|$ and $R_{\mathrm{reg}}(t)=-\lambda_r(r_t-r_{\mathrm{init}})^2$, where $\lambda_s=\lambda_r=0.1$ during training. Both terms are removed in the reward ablation, while the radius remains in the observation and action spaces. The terms discourage abrupt radius changes and drift from the initial radius; neither prescribes a target radius based on speed or obstacle distance. Our interpretation is that, without these terms, unstable radius choices can make future observations less consistent, especially at higher speeds. The ablation removes both terms together, so it does not isolate either term's causal contribution.
\Changes
\begin{itemize}
    \item Sec.~IV-B-3 adds the voxel resolution--coverage trade-off, the speed-dependent performance gap, and the observed abrupt deceleration and loss of balance, together with the limited-look-ahead explanation.
    \item The Active Perception item in Sec.~III-B-3 now includes the formulas for the radius smoothness and initial-radius deviation penalties and defines the radius variables.
    \item Sec.~IV-B-4 specifies that the reward ablation removes both terms while retaining the radius in the observation and action spaces, and discusses their role in stabilizing radius selection at higher speeds (Table~V).
\end{itemize}
\Location{Sec.~III-B-3, p.~4; Secs.~IV-B-3--IV-B-4, pp.~5--6; Tables~III and V, p.~6.}



\subsection*{AE Comment 3: Clarifications and presentation}
\begin{reviewcomment}
Resolve whether the radius changes the vertical map only or both maps. State whether the adaptive-radius ablation was run in simulation or on hardware. Clarify whether the robot recovers autonomously in the drop experiment or the safety tether arrests the fall. Add confidence intervals or variance to Fig. 5, and correct the likely typo in the Introduction.
\end{reviewcomment}
\Response Thank you for these constructive comments. We agree that these details should be clarified. The radius affects only the vertical distance map; horizontal-map coverage and resolution are fixed. The Table~V adaptive-radius ablation was run in simulation. In the repeated drop experiment, the robot recovered autonomously; the tether was a hardware safety measure and did not arrest the fall. We reduced its tension and repeated the experiment ten times in the same scene with the same policy to clarify its role. We have replaced the drop-experiment video clip in the Multimedia Materials to make the role of the safety tether clearer. We revised Fig.~5 to show the mean and sample standard deviation of success rates across ten independently trained policies, each evaluated using its final checkpoint under the same protocol. We also removed the erroneous wording \textit{3-D negotiation tasks} from the Introduction.
\Changes
\begin{itemize}
    \item We clarified the maps in Sec.~I, Sec.~II-B, and Sec.~III-B-1.
    \item We identified the simulation ablation setting in Sec.~IV-B-4 and removed the Table~V comparison from Sec.~IV-C-1.
    \item We revised Fig.~5 and its caption in Sec.~IV-C-2.
    \item We corrected the typo in Sec.~I.
\end{itemize}


\Location{Sec.~I, p.~1; Sec.~II-B, p.~2; Sec.~III-B-1, p.~3; Sec.~IV-B-4, pp.~5--6; Table~V, p.~6; Sec.~IV-C-1, p.~7; Sec.~IV-C-2 and Fig.~5, p.~7. Tether clarification: response letter and updated drop-experiment video clip in the Multimedia Materials.}

\clearpage
\section*{Response to Reviewer 3}
\subsection*{Reviewer 3, General assessment}
\begin{reviewcomment}
This paper presents ADAPT, a perceptive locomotion framework for humanoid robots that decomposes local 3D geometry into two projections: a horizontal elevation map and a vertical distance map. The proposed approach is evaluated against several representative baselines and is further deployed on a Unitree G1 humanoid robot in the real world.
\end{reviewcomment}
\Response Thank you for your careful reading and for recognizing both the comparative evaluation and the real-world Unitree G1 deployment. We appreciate your constructive comments, which have helped us make the method, baseline comparisons, and deployment evidence clearer.

\subsection*{Reviewer 3, Comment 1: Motivation}
\begin{reviewcomment}
This paper addresses an important problem in perceptive locomotion, and the motivation is well established.
\end{reviewcomment}
\Response Thank you for recognizing the importance and motivation of this work. ADAPT aims to coordinate limited perceptual resolution with the different look-ahead requirements of locomotion in complex three-dimensional environments.


\subsection*{Reviewer 3, Comment 2: Adaptive perception radius}
\begin{reviewcomment}
Incorporating the perception radius into the policy's action space is a clever design choice. It allows the robot to adaptively adjust its perceptual range according to the requirements of different 3D locomotion tasks.
\end{reviewcomment}
\Response Thank you for this assessment. Including the radius in the action space lets the policy learn locomotion and sensing-range adjustment together. The vertical map's input dimension stays fixed while its radial coverage and sampling interval vary.


\subsection*{Reviewer 3, Comment 3: Voxel-based baseline analysis}
\begin{reviewcomment}
Table III shows that ADAPT outperforms all baselines, including some computationally intensive voxel-based approaches. Given that voxel representations retain richer 3D geometric information, it would be helpful for the authors to provide further analysis explaining why these methods perform worse than ADAPT. In particular, the paper should clarify whether the performance gap arises from optimization difficulty, limited network capacity, training-budget constraints, representation redundancy, or other implementation factors.
\end{reviewcomment}
\Response Thank you for this suggestion. We have added an analysis in Sec.~IV-B-3 (Performance Comparison). Voxel representations retain richer 3D geometry and are not inherently inferior. For the Gallant voxel baseline implemented in this study, however, fine spatial resolution limits sensing coverage under a fixed voxel count and computational budget. The success-rate gap relative to ADAPT is smaller at lower speeds and widens at high speed. In high-speed tests, we observed abrupt deceleration after obstacles entered the sensing range, sometimes followed by loss of balance. We attribute this behavior to limited look-ahead, which may leave insufficient time for smooth adjustments of speed and posture and thus increase the difficulty of high-speed traversal. In contrast, ADAPT allows the policy to adjust its sensing range, and hence its look-ahead distance, according to the robot's motion state.
\Changes
\begin{itemize}
    \item Sec.~IV-B-3 now explains that, for the implemented Gallant baseline, fine voxel resolution restricts sensing coverage under a fixed voxel count and computational budget.
    \item The same section describes the smaller performance gap at lower speeds and the abrupt deceleration and occasional loss of balance observed at high speed, explaining how limited look-ahead can hinder smooth speed and posture adjustments.
    \item The analysis acknowledges that voxels preserve richer 3D geometry and contrasts the implemented Gallant baseline's fixed sensing coverage with ADAPT's policy-controlled look-ahead distance.
\end{itemize}
\Location{Sec.~IV-B-3, p.~5; Table~III, p.~6.}



\subsection*{Reviewer 3, Comment 4: Drop experiment and tether}
\begin{reviewcomment}
The real-world experiments provide valuable hardware validation. However, for the drop experiment, it is difficult to determine whether the robot recovers its balance autonomously or whether its fall is arrested by the safety tether.
\end{reviewcomment}
\Response Thank you for pointing out the video ambiguity. In our repeated drop experiment, the robot recovered autonomously; the tether served as a hardware safety measure rather than arresting the fall. We reduced its tension and repeated the experiment ten times in the same scene with the same policy to clarify its role. We have replaced the drop-experiment video clip in the Multimedia Materials to make the role of the safety tether clearer.
\Changes We have replaced the drop-experiment video clip in the Multimedia Materials.
\Location{Response letter and Multimedia Materials (updated drop-experiment video clip); no manuscript text change.}



\subsection*{Reviewer 3, Comment 5: Deployment latency and onboard resources}
\begin{reviewcomment}
Table IV reports training efficiency, but inference efficiency is arguably more important for perceptive locomotion because training is performed offline, whereas perception and control must operate under strict real-time constraints during deployment. The authors are encouraged to report inference-time benchmarks, including map-construction latency, policy inference latency, total perception-to-control latency, and onboard computational resource usage. A comparison of these quantities across ADAPT and the main baselines would better support the paper's efficiency claims.
\end{reviewcomment}
\Response Thank you for this suggestion. We have extended Table~IV with runtime measurements for ADAPT and the Gallant voxel baseline and added the measurement setup and analysis in Sec.~IV-C-3. Both methods were deployed and benchmarked on the same local computer with an Intel i7-12700H CPU and an NVIDIA RTX 3060 Laptop GPU with 6 GiB VRAM. Both the real-world robot experiments and the latency and resource-usage benchmarks used this local-computer deployment setup. Perception and policy inference ran on the local computer, which generated and sent joint-angle targets to the robot at 50 Hz. No Jetson Orin NX was used in either set of experiments. We have corrected the original deployment description in Sec.~IV-A-2 and clarified the shared setup in Sec.~IV-C-3. Table~IV(b) reports latency and resource usage on the local computer used for deployment.

Each method was tested five times on each of three terrain categories: flat ground, a single obstacle, and a complex environment, for 15 trials per method. Latency entries are equal-weight averages of per-trial means and P95s. Map latency spans the input callback to the first observed map-ready event; perception-to-publish latency spans that callback to the first command-publication call using the update. Policy inference latency on GPU is measured with current-stream CUDA events. For ADAPT, each map-latency and perception-to-publish entry reports the larger statistic of its two map streams. CPU utilization and RAM usage are whole-host averages, and peak VRAM is the maximum sampled usage.

Relative to Gallant, ADAPT reduces mean map latency from 64.06 to 41.64 ms and mean perception-to-publish latency from 79.45 to 55.94 ms. Mean policy inference latency on GPU decreases from 1.91 to 1.40 ms, and sampled peak VRAM decreases from 5,711 to 2,373 MiB. Mean CPU utilization decreases from 38.19\% to 17.76\%, and RAM usage from 10.66 to 8.04 GiB.
\Changes Sec.~IV-A-2 now states that real-world perception and policy inference run on the local computer, which sends joint-angle targets to the robot at 50 Hz. Table~IV(b) and Sec.~IV-C-3 report the runtime comparison on the same deployment computer, measurement boundaries, hardware, and trial setup.


\Location{Abstract and Sec.~I, p.~1; Sec.~IV-A-2, p.~4; Table~IV(b), p.~6; Sec.~IV-C-3, p.~7.}

\clearpage
\section*{Response to Reviewer 8}
\subsection*{Reviewer 8, General assessment}
\begin{reviewcomment}
The paper proposes ADAPT, a perceptive locomotion framework for humanoid robots. ADAPT decomposes local 3D geometry into two compact projections: a horizontal elevation map for footholds and a vertical distance map for body-level clearance. Moreover, it adds the sensing radius as an extra action of the reinforcement learning policy. The authors compare ADAPT in simulation against PIM, Humanoid Parkour, a TSDF, and Gallant, and deployed zero-shot on a Unitree G1.

To my knowledge, the content of the article is novel and sufficiently significant. The authors justify their method well and are clear in the explanation of the limits of baseline methods and the strengths of their proposed technique. The experiments validate the authors' assumptions. The organization of the paper is adequate, with one exception listed below. To my knowledge, the citations are adequate. I believe this to be a valuable contribution, still I believe some parts could be clearer.
\end{reviewcomment}
\Response We are grateful for your encouraging assessment of the novelty and significance of ADAPT and for your recognition of the motivation, comparisons, and experimental validation. Thank you also for noting that the organization and citations are generally adequate. Your detailed suggestions have helped us resolve the remaining presentation issue and explain the method and results more precisely.

\subsection*{Reviewer 8, Comment 1: Typo}
\begin{reviewcomment}
There is a possible typo in the introduction: 3-D negotiation tasks, you mean 3-D navigation tasks?
\end{reviewcomment}
\Response Thank you for catching this error. We agree and have corrected it. We removed the erroneous wording \textit{3-D negotiation tasks} when revising Sec.~I (Introduction).
\Location{Sec.~I, p.~1.}



\subsection*{Reviewer 8, Comment 2: Support for the efficiency claims}
\begin{reviewcomment}
In the introduction, you claim that the adaptive sensing mechanism is responsible for reducing the computational overhead by 56\%. Where in the experiment section do you support this claim? The same question holds for the 98.6\% reduction in observation dimensionality with respect to voxel-based methods that is claimed in the abstract. I think this comparison could be more clearly stated.
\end{reviewcomment}
\Response Thank you for pointing this out. We agree that the original comparison was imprecise. The 56\% result is the reduction in training \emph{peak GPU memory} relative to Gallant: $(36{,}054-15{,}870)/36{,}054=55.98\%$, with values in MiB. Using the same modified NVIDIA RTX 4090 GPU with 48 GB of VRAM and common PPO settings (4,096 parallel environments, a rollout of 24 steps per environment per policy update, 5 PPO epochs per update, and 4 mini-batches per epoch; each rollout batch contains $4{,}096\times24=98{,}304$ transitions), iteration time is 34.6\,s for Gallant and 11.3\,s for ADAPT; ADAPT takes 32.7\% of the baseline's time. These results compare the implemented representation and encoder pipelines, not the isolated effect of adaptive-radius selection.

Sec.~IV-B-3 states that the methods share the same training settings; the specific numerical PPO settings listed above are provided only in this response letter.

The 98.6\% result compares only single-frame perception inputs \emph{before encoding}: $32\times32\times40=40{,}960$ voxel values for Gallant versus $21\times17+13\times17=578$ map values for ADAPT, giving $(40{,}960-578)/40{,}960=98.59\%$. It does not compare total actor observations.
\Changes
\begin{itemize}
    \item The Abstract states a 98.6\% reduction in exteroceptive input dimensionality relative to our implementation of Gallant and qualitatively summarizes lower deployment latency and resource usage.
    \item Sec.~I gives the quantitative training and local-computer deployment comparisons.
    \item The caption of Table~I in Sec.~II-A defines sensory dimension as pre-encoder input elements, excluding proprioception, history, and critic-only inputs.
    \item Sec.~IV-B-3 gives the dimensionality calculation and states that hardware and training settings are identical across methods; Table~IV(a) reports wall-clock iteration times and peak GPU memory. The specific numerical PPO settings are provided only in this response letter.
\end{itemize}
\Location{Abstract and Sec.~I, p.~1; Sec.~II-A and Table~I, p.~2; Sec.~IV-B-3, p.~5; Table~IV(a)--(b), p.~6; Sec.~IV-C-3, p.~7.}



\subsection*{Reviewer 8, Comment 3: Which map the radius changes}
\begin{reviewcomment}
The abstract, introduction, and Sec. II-B say the radius modulates the vertical map only. Sec. III-B-1 says it acts as a dynamic gain modulating the field of view for both projection maps. Is the horizontal map also influenced by the radius?
\end{reviewcomment}
\Response Thank you for highlighting this ambiguity. We agree that the scope of radius adaptation should be explicit. The adaptive radius affects only the vertical distance map. Horizontal-map coverage and sampling resolution remain fixed.
\Changes
\begin{itemize}
    \item Secs.~I and II-B clarify that radius adaptation applies only to the vertical distance map, while the horizontal elevation map retains fixed coverage and sampling resolution.
    \item Secs.~III-A and III-C identify the vertical-map sensing radius as a policy output; Sec.~III-B-1 identifies the previous radius $r_{t-1}$ as an observation.
    \item Sec.~III-B-2 explains how $r_t$ changes the vertical map's radial range and sampling interval, $\delta_r=r_t/N_r$, while the angular span and number of bins remain fixed. Fig.~2 illustrates this range--resolution trade-off.
\end{itemize}
\Location{Abstract and Sec.~I, p.~1; Secs.~II-B and III-A, p.~2; Secs.~III-B-1--III-B-2, p.~3; Sec.~III-C and Fig.~2, p.~4.}



\subsection*{Reviewer 8, Comment 4: Ablation setting}
\begin{reviewcomment}
Was the ablation study on the adaptive radius mechanism done in simulation or in the real-world? Discussing Table V results in Sec. IV-C-1 creates confusion.
\end{reviewcomment}
\Response Thank you for requesting this clarification. We agree that the evaluation setting should be explicit. The adaptive-radius ablation in Table~V was conducted in simulation.
\Changes
\begin{itemize}
    \item Sec.~IV-B-4 identifies the tests as simulation experiments, with ten trials per method, terrain, and speed command, and specifies the two removed radius-reward terms.
    \item The same section clarifies that each trial contains five obstacles and that success is the fraction of obstacles crossed, with statistics pooled across the three terrains at each speed.
    \item We removed the Table~V comparison from Sec.~IV-C-1, which covers hardware tests.
\end{itemize}
\Location{Sec.~IV-B-4, pp.~5--6; Table~V, p.~6; Sec.~IV-C-1, p.~7.}



\subsection*{Reviewer 8, Comment 5: Adaptive reward}
\begin{reviewcomment}
The ablation shows that without the adaptive reward the policy fails to learn the optimal radius as speed increases. Do you have an explanation on why that happens? Which part of the adaptive reward is responsible for retrieving the optimal policy? Is the optimal reward basically enforcing a hand-designed heuristic? Understanding this could help better define the scope of the contribution.
\end{reviewcomment}
\Response Thank you for this helpful comment. We agree that the reward formulation needs a clearer explanation. The training reward for adaptive perception contains a smoothness term and a radius regularizer. With current radius $r_t$, previous radius $r_{t-1}$, and initial radius $r_{\mathrm{init}}$, their contributions to the total reward can be written as
\[
R_{\mathrm{smooth}}(t)=-\lambda_s|r_t-r_{t-1}|,\qquad
R_{\mathrm{reg}}(t)=-\lambda_r(r_t-r_{\mathrm{init}})^2,
\]
where $\lambda_s=\lambda_r=0.1$ during training. Thus, the combined radius-related contribution to the total reward is $-0.1|r_t-r_{t-1}|-0.1(r_t-r_{\mathrm{init}})^2$. Both contributions are removed in the reward ablation, while the radius remains in the observation and action spaces and all other reward terms are retained. The terms penalize abrupt radius changes and drift from the initial radius. Neither sets a target radius from speed or obstacle distance, and training uses no target-radius tracking reward. The regularizer encodes a preference for the initial radius, while allowing task rewards to favor a different radius when beneficial; it does not prescribe a speed-to-radius rule.

Our interpretation is that radius actions affect locomotion indirectly through future perception, which may be harder to learn than direct motor actions. Without regularization, early radius oscillation or drift can make observations less consistent; higher speeds may increase sensitivity to unsuitable or unstable sensing ranges. This is an interpretation of the ablation, not a separate causal measurement of either reward term.
\Changes
\begin{itemize}
    \item The Active Perception item in Sec.~III-B-3 now includes the formulas for the smoothness and initial-radius deviation penalties and defines the radius variables.
    \item Sec.~IV-B-4 states that the reward ablation removes both terms and interprets the high-speed result.
    \item The formulas are reproduced above for convenience. The numerical training weights and the explicit clarification that no speed- or obstacle-dependent target radius is prescribed are provided only in this response letter.
\end{itemize}
\Location{Sec.~III-B-3, Active Perception item and reward formulas, p.~4; Sec.~IV-B-4, pp.~5--6; Table~V, p.~6. Numerical training weights and explicit target-radius clarification: response letter only.}



\subsection*{Reviewer 8, Comment 6: Actor inputs and recurrence}
\begin{reviewcomment}
Fig. 1 makes it look like the actor also receives the privileged states, which contradicts what is claimed by the text in Sec. III-C. I also think Fig. 1 could be more informative if it showed the recurrent nature of the actor.
\end{reviewcomment}
\Response Thank you for identifying the diagram error. We have revised Fig.~1 so that the actor receives only deployable observations while the separate critic receives the additional privileged elevation map. The figure now shows independent actor and critic LSTM hidden-state loops. Each branch uses an LSTM with a hidden dimension of 256, followed by an MLP with hidden dimensions $[128,64,32]$ and ELU activations. The horizontal-map MLP encoder has hidden dimensions $[256,128,64]$, and the vertical-map MLP encoder has hidden dimensions $[128,64]$; both use CELU activations and produce 32-dimensional embeddings. The two map embeddings form a 64-dimensional environmental representation. The actor and critic maintain separate recurrent states, and the privileged ground-truth elevation map is supplied only to the critic.
\Changes Fig.~1 in Sec.~III-B-1 (Observation Space) and Sec.~III-C (Training Architecture) now distinguish the inputs and recurrent states of the actor and critic.
\Location{Sec.~III-B-1 and Fig.~1, p.~3; Sec.~III-C, p.~4. Layer specifications: response letter.}



\subsection*{Reviewer 8, Comment 7: Uncertainty in Fig. 5}
\begin{reviewcomment}
Fig. 5 only shows the average success rates, it would be better to show also confidence intervals/variance. The plot as it is does not give information on how much performance varies across seeds.
\end{reviewcomment}
\Response Thank you for highlighting the variability across seeds. The experiments use ten policies independently trained with different random seeds, with the final checkpoint from each training run evaluated under the same protocol. For each terrain in each domain (simulation and real world), each policy is tested once at each of three command speeds ($0.5$, $1.0$, and $1.5$\,m/s), giving 30 binary outcomes. We first average the three outcomes for each policy, then report the mean and sample standard deviation of these ten policy-level success rates in Fig.~5. Thus, the error bars show variability across independently trained policies under this evaluation protocol; they are not confidence intervals computed by treating all 30 outcomes as independent. Because each policy is tested only once per speed, the measured variability can also reflect evaluation randomness. The simulation evaluation uses single-obstacle configurations matched to the fixed real-world scenes, distinct from the five-obstacle tracks in Table~III. Success means traversing the single obstacle.
\Changes
\begin{itemize}
    \item Sec.~IV-C-1 specifies ten independent training seeds, final-checkpoint selection, and the common evaluation protocol: one trial per policy at each of three speeds, for 30 trials per terrain and domain.
    \item Fig.~5 and its caption in Sec.~IV-C-2 now report the mean and sample standard deviation across the ten policy-level success rates and specify single-obstacle evaluation in matched simulation and real-world scenes.
\end{itemize}
\Location{Secs.~IV-C-1--IV-C-2, p.~7; Fig.~5, p.~7.}


\end{document}
